\hypersetup{pageanchor=false}
\chapter*{Kurzfassung}
\thispagestyle{empty}

Batterien sind eine Schlüsseltechnologie der Energiewende, da sie die Speicherung volatiler Energie aus erneuerbaren Quellen ermöglichen. Während die elektrochemische Energiedichte auf Zellebene kontinuierlich gesteigert wurde, erfordert der Mangel an strukturellen Eigenschaften bisher schwere Sicherheitsgehäuse. Diese passive Masse reduziert die effektive Energiedichte des Gesamtsystems erheblich. Das Konzept der Strukturbatterie, bei dem das Material gleichzeitig Lasten trägt und Energie speichert, adressiert dieses Problem und wurde vom World Economic Forum als eine der "`Top 10 Emerging Technologies"' ausgezeichnet.

Um die Identifizierung leistungsfähiger Materialkombinationen zu beschleunigen, wird in dieser Arbeit ein multiskaliger und multiphysikalischer Ansatz zur Untersuchung von Strukturbatterien vorgestellt. Ausgehend von einer Analyse der Einzelkomponenten wird eine Simulationsumgebung entwickelt, die das Verhalten von der Mikro- bis zur Makroebene abbildet. Während die Gesamtsimulation maßgeblich auf der Finite-Elemente-Methode basiert, wurden für den Strukturelektrolyten spezifische Modellierungsansätze, wie etwa als Poren-Netzwerk oder mittels Walk-On-Stars, implementiert, um die effektiven Eigenschaften mit minimalem Rechenaufwand zu ermitteln.

Darauf aufbauend wurde ein analytischer Berechnungsansatz abgeleitet, der eine schnelle Vorhersage der Energiedichte sowie der Zug- und Biegesteifigkeit erlaubt. Im Rahmen dieser Arbeit wurde dieser Ansatz genutzt, um erstmalig eine systematische Untersuchung von 6.600 Materialkombinationen durchzuführen. Die Ergebnisse identifizieren mehrere hochgradig vielversprechende Konfigurationen und liefern damit eine essenzielle Grundlage für die zukünftige Materialentwicklung im Bereich der multifunktionalen Energiespeicher.

\pagebreak
\thispagestyle{empty}
{\bfseries \LARGE Abstract}
\vspace{1cm}

Batteries are a key technology in the energy transition, as they enable the storage of volatile energy from renewable sources. While electrochemical energy density at the cell level has been continuously increased, the lack of structural properties has so far required heavy safety housings. This passive mass significantly reduces the effective energy density of the overall system. The concept of the structural battery, in which the material simultaneously carries loads and stores energy, addresses this problem and has been recognized by the World Economic Forum as one of the “Top 10 Emerging Technologies.”

To accelerate the identification of high-performance material combinations, this work presents a multiscale and multiphysical approach to the investigation of structural batteries. Based on an analysis of the individual components, a simulation environment is developed that maps the behavior from the micro to the macro level. While the overall simulation is largely based on the finite element method, specific modeling approaches, such as pore networks or walk-on stars, were implemented for the structural electrolyte in order to determine the effective properties with minimal computational effort.

Based on this, an analytical calculation approach was derived that allows rapid prediction of energy density as well as tensile and flexural stiffness. As part of this work, this approach was used to conduct a systematic investigation of 6,600 material combinations for the first time. The results identify several highly promising configurations and thus provide an essential basis for future material development in the field of multifunctional energy storage devices.
\cleardoublepage
\hypersetup{pageanchor=true}